\chapter{INTRODUÇÃO}
\label{introducao}

% \textcolor{red}{(Importante: somente a partir da introdução as páginas do trabalho são numeradas. Utilizam-se
% algarismos arábicos, sendo que, a contagem das páginas inicia na folha de rosto.)}

% O texto deverá ser digitado em espaço 1,5. O parágrafo deverá apresentar um recuo na primeira linha a 1,25cm da margem esquerda, não contendo espaçamento entre um parágrafo e outro. 

% Devem ser digitados em espaços simples: as citações com mais de 3 linhas (citações longas), notas, resumo, referências, legendas de ilustração e de tabelas, bem como partes da capa e da folha de rosto. 

% Os títulos das seções devem ser separados do início do texto que os precedem ou os sucedem por um espaço (1,5).

Existem centenas de milhões de investidores no mundo e eles podem ser classificados majoritariamente em três tipos, de acordo com \citeonline{investir_perfil_2022}: conservador, moderado e arrojado (ou agressivo). O conservador prioriza a preservação do capital e a alta liquidez; este perfil possui baixa tolerância à volatilidade e prefere maior previsibilidade e estabilidade. Já o moderado busca um equilíbrio entre segurança e rentabilidade. E o arrojado foca na maximização do patrimônio a longo prazo, possuindo alta tolerância a oscilações bruscas e perdas temporárias, compreendendo que a exposição a ativos de maior risco é necessária para obter retornos significativamente acima da média do mercado.

Índices de mercado, como o S\&P 500, são carteiras teóricas de ativos que funcionam como indicadores de referência, cujo objetivo é refletir o comportamento médio e a evolução dos preços de um determinado segmento da economia, como tecnologia e industria por exemplo \cite{assaf_neto_mercado_2022}. Os índices mais usados e famosos referentes ao mercado de ações norte-americano, como NASDAQ e \textit{Dow Jones Industrial Average} tendem a concentrar risco em empresas supervalorizadas; isso ocorre porque os preços das ações são compostos pelo valor mais um ruído, ou seja, divergências temporárias entre o valor de mercado e o valor intrínseco do ativo. Como o índice ponderado por capitalização usa o preço para definir o peso, ele matematicamente atribui mais peso a empresas onde o ruído é positivo (sobreavaliação) e menos peso onde o ruído é negativo (subavaliação) \cite{siegel_future_2005}.

Decorre dessa característica o problema central desta pesquisa: os índices de mercado consolidados são instrumentos deliberadamente generalistas, concebidos para representar o comportamento médio de um segmento e não para atender às necessidades específicas de cada perfil de investidor. Um mesmo índice é oferecido indistintamente ao conservador e ao arrojado, ainda que suas tolerâncias ao risco sejam opostas. Abre-se, assim, espaço para o estudo de índices alternativos, construídos a partir de análises históricas do mercado de ações e centrados no risco do ativo, e não estritamente em seu preço ou valor de mercado, de modo a comportar qualquer tipo de investidor.

A fundamentação desta pesquisa ancora-se na premissa de que a alocação de ativos é condicionada à tolerância ao risco do investidor. Segundo \citeonline{bodie_essentials_2024}, a categorização dos investidores em perfis não é meramente descritiva, mas define os parâmetros de restrição em modelos de otimização de portfólio. A relevância do problema, portanto, não é apenas técnica, mas prática: uma análise que ignore o risco e o perfil entrega ao investidor uma recomendação descalibrada para seus próprios objetivos. Nesse contexto, a análise de métricas técnicas serve como mecanismo de filtragem capaz de explorar ineficiências de mercado para atender expectativas distintas: maximizar o retorno no longo prazo, reduzir a volatilidade, aumentar a previsibilidade, entre outras.

A viabilidade dessa proposta depende primeiro da qualidade do dado. Segundo \citeonline{PALESTINO2014}, em ambientes de \textit{Business Intelligence}, o tratamento do dado é o estágio mais crítico, pois é onde se aplicam a limpeza e a padronização necessárias para eliminar inconsistências, e sem essa etapa qualquer análise subsequente herda as inconsistências da entrada. Além disso, a qualidade das análises também depende da validação sobre dados históricos, um caminho objetivo para examinar como uma regra de seleção de ativos teria se comportado sob condições reais de mercado, dentro dos limites das condições de validação adotadas em cada aplicação.

\section{OBJETIVO GERAL}
\label{sec:objetivo-geral}

O objetivo geral deste trabalho é construir e avaliar um ambiente analítico integrado para a análise de risco e retorno de ações do mercado norte-americano, capaz de tratar os dados, popular um \textit{Data Mart} baseado em um modelo dimensional e viabilizar a criação de novos índices de mercado sob a perspectiva de diferentes perfis de investidor. Esse ambiente compreende, como componentes, um fluxo de \textit{Extract, Transform and Load} (ETL) para a coleta, o tratamento e o armazenamento sistemático de dados acionários no período de 2010 a 2025, o cálculo de uma métrica própria de risco e um mecanismo de simulação histórica de carteiras, validados em comparação ao índice de referência S\&P 500.

\section{Objetivos Específicos}
\label{sec:objetivos-especificos}

Para atingir o objetivo geral, este trabalho se organiza em torno de objetivos específicos:

\begin{itemize}
    \item Identificar as fontes de dados acionários e institucionais a serem adotadas, delimitando o universo de ativos considerados e o índice utilizado como referência de mercado.
    \item Elaborar a modelagem dimensional do \textit{Data Mart}, definindo tabelas fato e dimensão que sustentem as consultas analíticas posteriores.
    \item Operacionalizar o processo de transformação sobre essa modelagem, especificando as etapas de limpeza, padronização e carga que garantem a consistência dos dados carregados.
    \item Formalizar a métrica composta de risco e seus modificadores.
    \item Confrontar o modelo com o índice de referência sob a ótica das cinco faixas de risco derivadas da métrica composta, examinando o comportamento comparativo entre carteiras selecionadas pelo modelo e o \textit{benchmark} ao longo do intervalo de análise.
\end{itemize}

\section{Estrutura do Trabalho}

O restante deste documento está organizado da seguinte forma: O Capítulo~\ref{fundamentacao} apresenta a fundamentação teórica, abordando os perfis de investidor, as métricas quantitativas de risco, os indicadores empregados como modificadores, a justificativa da modelagem dimensional e a revisão de literatura dos trabalhos relacionados. O Capítulo~\ref{metodologia} descreve a metodologia em nível conceitual, detalhando a coleta, o tratamento, a modelagem do grau de risco e a estratégia de validação. O Capítulo~\ref{implementacao} relata a implementação prática, incluindo a \textit{stack} tecnológica, a arquitetura do \textit{Data Mart}, o \textit{pipeline} de ETL e a construção do \textit{dashboard}. O Capítulo~\ref{chap:experimentacao} apresenta a experimentação e a avaliação comparativa dos resultados contra o S\&P 500. Por fim, o Capítulo~\ref{chap:conclusao} reúne as considerações finais, limitações e propostas de trabalhos futuros.

\bigskip

